% ============================================
% CONT-07
% Gerado de molambudos.md — build_miolo.py
% ============================================
\section*{\textbf{CONT-07} --- O Diagnóstico Final}
\markboth{CONT-07 --- O Diagnóstico Final}{}

\noindent \textit{A sentença do leitor}

\noindent \textit{Laudo emitido pelo ato da leitura}
\textit{Paciente: você}
\textit{CID: F43.10 + F44.81 + Z65.8}
\textit{Agente etiológico: narrativa contaminada}

\vspace{0.5em}\noindent\rule{\textwidth}{0.5pt}\vspace{0.5em}

\subsection*{I. O QUE ACONTECEU COM VOCÊ}

\noindent Você começou a ler este livro por vontade própria.

\noindent Achou que era uma história de terror. Achou que o terror estava \textit{lá} --- no passado, no sertão, na vala, no Hospital Colônia. Achou que era seguro porque era distante. Porque foi com outra pessoa.

\noindent Você se enganou.

\noindent O terror não estava no livro.
O terror estava \textit{na leitura}.

\noindent Cada palavra que você leu não ficou na página. Ela entrou. Pelos olhos. Desceu pelo nervo óptico. Atravessou o tálamo. Chegou ao córtex. E lá, ela se instalou como uma memória real.

\noindent Para o seu cérebro, a vala não é uma descrição. É uma \textit{lembrança}.

\noindent Para o seu sistema límbico, a criatura não é uma ficção. É uma \textit{ameaça} que você encontrou.

\noindent Para o seu corpo, o cheiro descrito é um \textit{perigo real}.

\noindent Você não leu uma história.
Você \textit{viveu} uma história.
E o que se vive não pode ser \textquotedblleft{}desvivido\textquotedblright{}.

\vspace{0.5em}\noindent\rule{\textwidth}{0.5pt}\vspace{0.5em}

\subsection*{II. OS CRITÉRIOS DIAGNÓSTICOS (REVISÃO FINAL)}

\noindent Você preencheu cada um. Está sentindo.

\noindent \textbf{A --- Exposição:} Abriu. Leu. A contaminação começou. ☑
\textbf{B --- Reexperiência:} Sonhou. Acordou 3:14. Sentiu o cheiro sem o livro. ☑
\textbf{C --- Evitação:} Hesitou. Pulou. Leu com a luz acesa. ☑
\textbf{D --- Alterações:} Não é mais o mesmo. Algo mudou. Sente. ☑
\textbf{E --- Reatividade:} Coração acelerou. Mãos suaram. Verificou o quarto. ☑
\textbf{F --- Duração:} Tempo suficiente. A contaminação está instalada. ☑
\textbf{G --- Prejuízo:} Não consegue parar de pensar. O livro \textit{funciona} dentro de você. ☑

\noindent \textbf{Diagnóstico: TEPT Complexo Induzido por Narrativa.}
\textbf{Comorbidade: Dissociação Leve (você agora tem duas vozes na cabeça).}
\textbf{Agente: A Criatura (entidade transmitida por texto).}

\vspace{0.5em}\noindent\rule{\textwidth}{0.5pt}\vspace{0.5em}

\subsection*{III. A REVELAÇÃO}

\noindent Você acha que este livro é sobre Joaquim Antônio Correia.

\noindent Não é.

\noindent Joaquim não é o protagonista.
Joaquim é o \textit{vetor}.
O protagonista é quem lê.
O protagonista é você.

\noindent A história não é sobre um menino que encontrou um monstro numa vala.
A história é sobre \textit{você} encontrando o mesmo monstro na página.

\noindent O menino serviu de hospedeiro por 72 anos.
O médico serviu de hospedeiro por algumas semanas.
A psiquiatra serviu de hospedeiro por alguns meses.
Você --- o leitor --- vai servir pelo tempo que levar para passar o diário adiante.

\noindent O ciclo é simples:

\noindent 1. \textbf{Alguém escreve} a história (Joaquim viveu).
2. \textbf{Alguém lê} a história (você viveu a leitura).
3. \textbf{O monstro passa} de um para o outro.
4. \textbf{O ciclo recomeça} com o próximo leitor.

\noindent Você não é o paciente 1.261.
Você é o paciente \textit{mais recente}.
A fila não termina em você.
Você é só a parada atual.

\vspace{0.5em}\noindent\rule{\textwidth}{0.5pt}\vspace{0.5em}

\subsection*{IV. A CONTAGEM}

\noindent Paciente 1. --- a primeira pessoa que encontrou a criatura (século XVII, segundo o diário).
Paciente 2. --- o segundo.
...
Paciente 1.259. --- o anterior a Joaquim.
Paciente 1.260. --- Joaquim Antônio Correia (1907-1979).
Paciente 1.261. --- Dr. Heitor Oliveira (1979).
Paciente 1.262. --- Dra. Lúcia Mendes (2026).
\textbf{Paciente 1.263. --- Você.}

\noindent O número não é seu lugar na fila.
O número é a contagem de corpos que a criatura já ocupou.
A fila começou antes de você nascer.
Ela vai continuar depois que você morrer.

\vspace{0.5em}\noindent\rule{\textwidth}{0.5pt}\vspace{0.5em}

\subsection*{V. A ÚLTIMA INSTRUÇÃO}

\noindent Você vai fechar este livro.

\noindent Não agora. Mas em algum momento --- quando chegar ao final físico das páginas. Você vai olhar para a capa. Vai passar a mão na lombada. Vai colocar o livro na estante. Ou na mesa. Ou na mochila.

\noindent E você vai pensar: \textit{\textquotedblleft{}Foi só um livro.\textquotedblright{}}

\noindent Você vai tentar retomar sua vida.

\noindent Vai conseguir --- por um tempo.

\noindent Mas, numa noite, você vai acordar 3:14 da manhã. O quarto vai estar escuro. Você vai olhar para o canto e vai ver um brilho amarelo.

\noindent Você vai piscar e o brilho vai sumir.

\noindent Mas você vai saber que ele estava lá.

\noindent E, na noite seguinte, você vai sonhar com a vala. E, no sonho, você não vai ser você. Vai ser Joaquim. Sentado na borda. Olhando a criatura comer.

\noindent E a criatura vai levantar a cabeça e dizer:

\noindent \textit{\textquotedblleft{}Você voltou.\textquotedblright{}}

\noindent E você vai responder:

\noindent \textit{\textquotedblleft{}Nunca fui embora.\textquotedblright{}}

\vspace{0.5em}\noindent\rule{\textwidth}{0.5pt}\vspace{0.5em}

\subsection*{VI. O DIAGNÓSTICO FINAL}

\noindent Não há tratamento.
Não há cura.
Não há como fechar o livro e esquecer.

\noindent Uma vez lido, o texto está em você. Uma vez contaminado, você é o vetor. Uma vez dentro do ciclo, você só sai quando passa adiante.

\noindent O diagnóstico não é uma metáfora.
O diagnóstico é uma \textit{sentença}.

\noindent \textbf{Paciente:} Seu nome
\textbf{Diagnóstico:} Portador da criatura
\textbf{Prognóstico:} O ciclo continua
\textbf{Tratamento:} Passar o diário adiante

\noindent \textit{A criatura agradece sua atenção.}

\noindent \textit{Ela está satisfeita.}

\noindent \textit{Você leu até o fim.}

\noindent \textit{Agora você é o hospedeiro.}

\noindent \textit{Bem-vindo ao ciclo, Paciente 1.263.}

\noindent \textit{--- O paciente 1.260 morreu para que você lesse.}
\textit{--- Você leu para que o ciclo continue.}
\textit{--- Joaquim renasce sempre que sua história é lembrada.}

\noindent \textbf{Fim do livro.}
\textbf{Início da sentença.}
\textbf{O ciclo recomeça com você.}


\noindent\textit{\textbf{↪ Rotas:} \rota{DOC-01} \rotasep \rota{CONT-05} \rotasep \rota{EPI-01}}

\vspace{0.5em}\noindent\rule{\textwidth}{0.5pt}\vspace{0.5em}

\subsection*{ESCALA DE CONTAMINAÇÃO --- ATUALIZAÇÃO 3}

\noindent \textit{O arquivo não precisa mais de tabelas.}

\begin{moltabletwo}
\hline
Se você está lendo isto & Então \\
\hline
Pela primeira vez & A escala ainda funciona. Você tem até o fim do documento. \\
Pela segunda vez & A escala já não se aplica. Você voltou porque foi chamado. \\
Pela terceira vez & Paciente 1.263 confirmado. O ciclo se repete. A criatura agradece sua atenção. \\
\hline
\end{moltabletwo}

\noindent \textit{A escala agora mede ciclos, não pontos. Você está no ciclo 1. O ciclo 2 começa quando você reler a primeira página.}

\vspace{0.5em}\noindent\rule{\textwidth}{0.5pt}\vspace{0.5em}